\begin{recipe}{Hartige tompouce met geitenkaasmousse}
\kicker[Voorgerecht,Kerst]{Voorgerecht · Kerst}
\index[register]{Hartige tompouce}
\index[register]{Geitenkaas}
\index[register]{Bladerdeeg}
\index[register]{Prosciutto}

\meta{30 minuten \dvd Oven 225\,°C}

\lettrine{E}{en} voorgerecht dat er feestelijk uitziet maar weinig werk vraagt.
De mousse kun je van tevoren maken en koud bewaren, zodat je vlak voor het eten
alleen nog hoeft op te maken.

\blockrule{Ingrediënten}
\begin{ingredients}
  \ing{75 g}{zachte geitenkaas}
  \ing{80 g}{prosciutto}
  \ing{5 plakjes}{bladerdeeg}
  \ing{1}{ei}
  \ing{100 ml}{slagroom}
  \ing{40 g}{rucola}
  \ing{40 g}{zongedroogde tomaatjes}
  \ing{20 g}{maanzaad}
  \ingb{balsamicocrema}
  \ingb{versgemalen peper}
\end{ingredients}

\blockrule{Bereiding}
\tip{Een dikkere balsamicoazijn werkt beter dan gewone balsamicocrema: hij blijft mooi liggen.}
\begin{steps}
  \item Verwarm de oven voor op 225 graden (boven- en onderwarmte).
  \item Halveer de plakjes bladerdeeg en leg ze op een met bakpapier beklede
    bakplaat. 
  \item Prik elk plakje een aantal keer in met een vork, dat voorkomt dat het deeg te
    veel opbolt. Bestrijk ze met losgeklopt ei.
  \item Verdeel het maanzaad over de helft van de plakjes: dit worden
    de bovenkantjes. 
  \item Bak alles in 10--12 minuten goudbruin en laat afkoelen.
  \item Roer de geitenkaas glad met een spatel. Klop in een aparte kom de
    slagroom stijf. 
  \item Spatel de slagroom voorzichtig door de geitenkaas en breng
    op smaak met zwarte peper. 
  \item Doe de mousse in een spuitzak en laat opstijven
    in de koelkast.
  \item Spuit de geitenkaasmousse in een zigzagbeweging op de plakjes bladerdeeg
    zonder maanzaad. 
  \item Leg per tompouce 2 plakjes prosciutto op de mousse en
    verdeel daar de rucola en zongedroogde tomaatjes over. Leg het
    maanzaadplakje er als deksel op.
  \item Druppel er wat balsamicocrema over en serveer meteen.
\end{steps}
\heroimagefade{images/tompouce-geitenkaas}
\end{recipe}
